% 第1章 为什么要学习编程（教材第1章 1.1--1.5）
\lstset{language=C}

\section{学习编程的热潮}

\subsection{席卷全球的“编程一小时”}

2013 年，公益网站 code.org 发起了每年 12 月持续一周的活动 Hour of Code，中文译作“编程一小时”。它的定位是“介绍计算机科学一小时的活动，旨在揭秘编程并说明它是任何人都可以学习的基础知识”。这项活动得到时任美国总统奥巴马以及扎克伯格、比尔·盖茨等科技界人士的支持，迅速从美国扩散到全球。

这件事说明的并不是“人人要当程序员”，而是：编程已经由专业技能变成了信息世界的通用工具，和识字、算术一样属于基础素养。

\subsection{资本汹涌的少儿编程}

2017--2018 年，资本大量进入少儿编程赛道，逻辑并不是“程序员缺口大”，而是押注“未来受教育人口普遍要接受编程教育”。资本判断的是趋势，不是课程内容本身——看懂这一点，就不必被“风口”牵着走。

\subsection{“再不学编程就晚了”}

这种恐慌来自 2016 年 AlphaGo 以 4:1 战胜李世石。但需要注意的是：AlphaGo 靠的是\textbf{高速搜索加海量数据存储}，并不是像人一样“思考”。它写不出新的算法思想，只能在你给定的规则空间里搜索。所以“学会编程就能凌驾于人工智能之上”并不成立，真正稀缺的是\textbf{提出问题和设计求解方法}的能力。

\corepoint{把“学编程”理解成“学一门手艺”会越学越窄；理解成“学一种描述计算过程的语言”才会越学越宽。前者随工具淘汰，后者可以迁移到任何语言。}

\section{为什么要学编程}

“人人学编程”不等于“人人都是程序员”。程序员不是人人都想当，也不是人人都能做。学编程的普遍意义在于获得\textbf{计算思维}（Computational Thinking）。

这一概念由美国卡内基·梅隆大学计算机科学系前系主任周以真（Jeannette Wing）于 2006 年在《Computational Thinking》一文中提出。她的核心论断是：计算思维代表着一种普遍的认识和一类普适的技能，每一个人——不仅仅是计算机科学家——都应热心于它的学习和运用。

把计算思维落到程序设计上，可以拆成四个可操作的要素：

\begin{table}[htbp]
\centering
\caption{计算思维的四个要素及其在 C 程序设计中的落点}
\small
\begin{tabularx}{\textwidth}{@{}lXl@{}}
\toprule
要素 & 含义 & 在 C 语言中的落点 \\
\midrule
抽象 & 丢掉与求解无关的细节，只保留决定结果的特征 & 用数据类型刻画出“一个整数”“一个字符”；用函数接口隐藏实现 \\
递归 & 把大问题化为同型的更小问题 & 递归函数必须有终止条件；栈帧的压入与弹出 \\
复用 & 一次写对，多处调用 & 库函数、自定义函数、头文件、宏 \\
折中 & 在正确、效率、可读、可移植之间取舍 & 用 double 保精度还是用 int 求速度；空间换时间 \\
\bottomrule
\end{tabularx}
\end{table}

教材给出的两条主线同样值得记住：一是\textbf{程序设计方法由自底向上走向自顶向下}，二是\textbf{数据结构由简单走向复杂}。前七章沿着“简单数据”这条线展开，第 8 章之后才进入数组、指针、字符串、结构体等复杂数据。

\section{什么是“编程”}

编程的实质是：把人对问题的求解思路，转写成计算机能够逐步执行、且每一步都无歧义的指令序列。这个转写过程包含三个层次：

\begin{itemize}
  \item \textbf{形式层}（语法）：写出来的东西必须符合语言的词法与语法规则，否则编译器拒绝翻译；
  \item \textbf{语义层}（含义）：每条语句在机器上真正做了什么——变量在内存的哪一格、值怎么变；
  \item \textbf{方法层}（思想）：把问题分解成什么样的步骤，用什么算法，数据结构怎么组织。
\end{itemize}

初学者最常见的困境是只盯着形式层：语法背得很熟，但拿到题不知道第一步写什么。原因就在于方法层没有建立——而这恰恰是这门课真正要训练的东西。

\corepoint{程序不是“写给人看的作文”，也不是“写给机器看的咒语”。它是\textbf{人与机器之间的合同}：人负责把方法想清楚，机器负责把每一步算到底。合同里任何一句有歧义，机器都会按它自己的理解执行，而且不会提醒你。}

\section{怎么学编程}

\subsection{一切都是计算}

计算机能做的事只有一件：按指令对内存里的数据进行极为简单的操作（读、写、加、比、跳），然后以极高的速度和极大的数量把简单操作堆成复杂结果。看清这一点，就不会再问“为什么计算机这么笨，连这个都不会”，而会问“这件事能不能拆成它会做的那些简单操作”。

\subsection{学编程不是学语法}

语法是必需的，但它只是入场券。教材的写法是从实际问题出发、在案例的不断深化中引出知识点；同一段程序会用四种写法改写（for / while / do-while / goto），目的不是让你背四种语法，而是让你看清\textbf{同一个算法思想可以用不同控制结构表达}。所以判断自己学会没有，不看“能不能背出语法格式”，而看“给定一个新问题，能不能拆出步骤并选对结构”。

\subsection{动手，动手，再动手}

\begin{enumerate}
  \item 每个知识点至少亲手敲一遍、编译一遍、跑一遍，不要只看。
  \item 故意制造错误：删掉初始化、写错分号、漏掉取地址符，看编译器报什么、运行结果怎么错。教材每章都有“常见错误小结”，正是为这一步准备的。
  \item 改需求：把“求 1 到 $n$ 之和”改成“求奇数之和”“求平方和”，检验自己是否真的掌握了模板而不只是记住了答案。
\end{enumerate}

\section{本章小结}

本章不含 C 语法，但它决定了后面 13 章的学习方式。三个结论：

\begin{itemize}
  \item 学编程的目的是获得计算思维（抽象、递归、复用、折中），不是背语法；
  \item 编程是把求解方法转写成无歧义指令序列的过程，难点在方法层；
  \item 唯一有效的路径是动手：敲、错、改、再改需求。
\end{itemize}

从第 2 章开始进入具体语法。届时请始终带着一个问题读代码：\textbf{这条语句在内存里到底改动了什么？}
